\documentclass[12pt,a4paper]{article}

% --- Paquetes Requeridos ---
\usepackage[utf8]{inputenc}
\usepackage[spanish]{babel}
\usepackage[T1]{fontenc}
\usepackage[margin=2.5cm]{geometry}
\usepackage{hyperref}
\usepackage{xcolor}
\usepackage{tabularx}
\usepackage{booktabs}
\usepackage{graphicx}
\usepackage{listings}
\usepackage{titlesec}
\usepackage{tocbibind}

% --- Configuración de Colores y Enlaces ---
\definecolor{primary}{RGB}{23, 54, 93}
\definecolor{secondary}{RGB}{47, 117, 181}
\definecolor{codegray}{rgb}{0.95,0.95,0.95}
\definecolor{codeblue}{rgb}{0.1,0.2,0.8}
\definecolor{codegreen}{rgb}{0.1,0.6,0.1}
\definecolor{codeorange}{rgb}{0.8,0.4,0.0}
\definecolor{errorred}{rgb}{0.8,0.1,0.1}

\hypersetup{
    colorlinks=true,
    linkcolor=primary,
    filecolor=secondary,      
    urlcolor=secondary,
    pdftitle={Manual Técnico V1 - MyAsistenciaBiometrical}
}

% --- Configuración de Bloques de Código ---
\lstset{
    backgroundcolor=\color{codegray},
    commentstyle=\color{codegreen},
    keywordstyle=\color{codeblue}\bfseries,
    stringstyle=\color{codeorange},
    basicstyle=\ttfamily\footnotesize,
    breakatwhitespace=false,         
    breaklines=true,                 
    captionpos=b,                    
    keepspaces=true,                 
    numbers=left,                    
    numbersep=5pt,                  
    showspaces=false,                
    showstringspaces=false,
    showtabs=false,                  
    tabsize=4,
    frame=single,
    rulecolor=\color{lightgray}
}

% --- Estilos de Títulos ---
\titleformat{\section}{\Large\bfseries\color{primary}}{\thesection}{1em}{}
\titleformat{\subsection}{\large\bfseries\color{secondary}}{\thesubsection}{1em}{}
\titleformat{\subsubsection}{\normalsize\bfseries\color{primary}}{\thesubsubsection}{1em}{}

% --- Inicio del Documento ---
\begin{document}

% --- PORTADA ---
\begin{titlepage}
    \centering
    \vspace*{4cm}
    
    {\Huge \bfseries \color{primary} MANUAL TÉCNICO V1 \par}
    \vspace{1cm}
    {\Large \bfseries Plataforma de Asistencia Biométrica Escolar y Gestión Académica \par}
    
    \vspace{4cm}
    
    \begin{flushleft}
        \large
        \textbf{Autor:} Marco Antonio Peña Vargas \\
        \textbf{Fecha:} Agosto 2026 \\
        \textbf{Versión:} 1.0.0 \\
        \textbf{Área:} Ingeniería en Sistemas Computacionales \\
        \textbf{Institución:} Tecnológico de Estudios Superiores del Oriente del Estado de México (TESOEM)
    \end{flushleft}
    
    \vfill
    {\large Documentación Oficial de Arquitectura, Configuración y Resolución de Errores\par}
\end{titlepage}

\tableofcontents
\newpage

% --- CONTROL DE VERSIONES ---
\section*{Control de Versiones}
\addcontentsline{toc}{section}{Control de Versiones}
\begin{table}[h]
    \centering
    \begin{tabularx}{\textwidth}{|c|c|X|X|}
        \hline
        \textbf{Versión} & \textbf{Fecha} & \textbf{Autor} & \textbf{Cambios} \\
        \hline
        1.0.0 & Agosto 2026 & Marco Antonio Peña Vargas & Definición de Arquitectura Limpia, Stack Tecnológico, Setup de Directorios, Configuración Global (Pydantic Settings) y Documentación de Errores de Inicialización. \\
        \hline
    \end{tabularx}
\end{table}

\newpage

% --- 1. ARQUITECTURA ---
\section{Arquitectura del Sistema}
El proyecto se desarrolla bajo el paradigma de \textbf{Clean Architecture} (Arquitectura Limpia) y un enfoque ligero de \textbf{Domain-Driven Design (DDD)}. Esta estructura garantiza que el código sea mantenible, testeable y preparado para entornos de producción, separando claramente las responsabilidades en capas (Presentación, Casos de Uso, Dominio e Infraestructura).

El flujo de información interno del Backend es estrictamente unidireccional y predecible:
\begin{center}
    \texttt{\textbf{Cliente $\rightarrow$ Router (API) $\rightarrow$ Service (Lógica) $\rightarrow$ Repository (Datos) $\rightarrow$ PostgreSQL (Supabase)}}
\end{center}

\subsection{Separación de Responsabilidades en el Backend}
\begin{itemize}
    \item \textbf{Routers:} Únicamente reciben la petición HTTP (ej. \texttt{POST /auth/login}), validan el esquema de entrada mediante \texttt{Pydantic} y delegan la ejecución. No contienen consultas SQL ni reglas de negocio.
    \item \textbf{Services:} Orquestan la lógica de negocio pura. Manejan el flujo de la aplicación (ej. Validar usuario $\rightarrow$ Validar contraseña $\rightarrow$ Crear JWT $\rightarrow$ Registrar auditoría).
    \item \textbf{Repositories:} Capa exclusiva para el acceso a datos (SELECT, INSERT, UPDATE, DELETE). Aisla al resto de la aplicación de la base de datos subyacente.
    \item \textbf{Domain / Models / Schemas:} Entidades que representan el modelo del negocio y validadores estrictos de entrada/salida.
\end{itemize}

% --- 2. STACK TECNOLÓGICO ---
\section{Stack Tecnológico y Librerías}

\subsection{Backend (Python 3.12)}
\begin{table}[h]
    \centering
    \begin{tabularx}{\textwidth}{|l|X|}
        \hline
        \textbf{Librería / Tecnología} & \textbf{Propósito dentro del Sistema} \\
        \hline
        \textbf{FastAPI} & Framework asíncrono para la creación de la API REST de alto rendimiento. \\
        \hline
        \textbf{Uvicorn} & Servidor web ASGI para ejecutar la aplicación FastAPI. \\
        \hline
        \textbf{Pydantic v2} & Validación estricta de datos (DTOs) y gestión de variables de entorno mediante \texttt{pydantic-settings}. \\
        \hline
        \textbf{SQLAlchemy} & ORM utilizado para interactuar con la base de datos relacional orientada a objetos. \\
        \hline
        \textbf{asyncpg} & Driver asíncrono optimizado para conectar con PostgreSQL (Supabase). \\
        \hline
        \textbf{Alembic} & Herramienta de migración de base de datos vinculada a SQLAlchemy. \\
        \hline
        \textbf{python-jose} & Generación, firma y validación criptográfica de tokens JWT. \\
        \hline
        \textbf{passlib (bcrypt)} & Hashing criptográfico seguro para almacenamiento de contraseñas. \\
        \hline
        \textbf{opencv-python} & Procesamiento de matrices de imágenes (\textit{headless} sin UI). \\
        \hline
        \textbf{face\_recognition} & Motor IA (dlib) encargado de generar embeddings faciales y compararlos. \\
        \hline
    \end{tabularx}
\end{table}

\subsection{Frontend (React Native \& Expo SDK 55+)}
\begin{table}[h]
    \centering
    \begin{tabularx}{\textwidth}{|l|X|}
        \hline
        \textbf{Librería / Tecnología} & \textbf{Propósito dentro del Sistema} \\
        \hline
        \textbf{Expo SDK} & Framework para compilar aplicaciones iOS/Android y web unificadas. \\
        \hline
        \textbf{React Navigation} & Enrutador principal para la navegación (Stacks, Tabs). \\
        \hline
        \textbf{Axios} & Cliente HTTP para el consumo de la API REST de FastAPI. \\
        \hline
        \textbf{expo-camera} & Acceso al hardware fotográfico para Kiosco y capturas de perfil. \\
        \hline
        \textbf{expo-location} & Acceso al GPS para validar el perímetro del salón de clases. \\
        \hline
        \textbf{expo-local-auth} & Interfaz con sensores biométricos nativos (FaceID/Huella dactilar). \\
        \hline
        \textbf{NativeWind} & Motor de estilos basado en la utilidad de Tailwind CSS. \\
        \hline
    \end{tabularx}
\end{table}

% --- 3. ESTRUCTURA DE CARPETAS ---
\section{Estructura de Directorios}
El proyecto adopta una arquitectura modular donde se aíslan las aplicaciones cliente, el servidor, la base de datos y el hardware físico.

\begin{lstlisting}[language=bash, caption=Estructura Raíz del Proyecto]
MyAsistenciaBiometrical/
|-- backend/                # API REST (FastAPI + Clean Architecture)
|   |-- .venv/              # Entorno virtual de Python
|   |-- app/                # Codigo fuente del backend
|   |   |-- api/            # Controladores y rutas HTTP
|   |   |-- application/    # Casos de uso
|   |   |-- core/           # Configuracion base, excepciones y middlewares
|   |   |-- database/       # Conexion, sesiones e inicializacion DB
|   |   |-- domain/         # Entidades puras y reglas de negocio
|   |   |-- infrastructure/ # Repositorios y servicios externos
|   |   |-- models/         # Clases de SQLAlchemy
|   |   |-- schemas/        # Modelos Pydantic (Request/Response)
|   |   |-- security/       # JWT y hashing
|   |   |-- services/       # Logica de orquestacion
|   |   `-- utils/          # Helpers y conversores
|   |-- tests/              # Pruebas unitarias y de integracion
|   |-- uploads/            # Archivos estaticos locales
|   |-- main.py             # Archivo de arranque Uvicorn
|   `-- requirements.txt
|-- Database/               # Scripts SQL (schema, triggers, views, RLS)
|-- docker/                 # Contenedores para BD local o pgAdmin
|-- docs/                   # UML, diagramas, manuales
|-- esp32/                  # Firmware C++ (Wi-Fi, HTTP, Relay)
`-- frontend/
    |-- kiosk/              # Proyecto Expo para Torniquete
    `-- mobile/             # Proyecto Expo principal (Alumnos/Maestros)
        `-- src/            # components, navigation, screens, hooks
\end{lstlisting}

% --- 4. CONFIGURACIÓN DEL BACKEND ---
\section{Configuración Base del Backend (\texttt{config.py})}

El archivo central de configuración de la plataforma se ubica en: \texttt{backend/app/core/config.py}. 
Su responsabilidad principal es gestionar, tipar y validar de forma estricta todas las variables de entorno necesarias para la ejecución segura del sistema, evitando que datos sensibles sean almacenados en texto plano en el código fuente.

\subsection{Funcionamiento de \texttt{pydantic-settings}}
La aplicación utiliza \texttt{BaseSettings} y \texttt{SettingsConfigDict} de la librería \texttt{pydantic-settings}. Esto proporciona:
\begin{enumerate}
    \item \textbf{Carga automática:} Busca variables en el sistema operativo o en el archivo \texttt{.env}.
    \item \textbf{Validación estricta de tipos:} Si una variable se define como \texttt{int} pero recibe un tipo diferente, FastAPI detiene el arranque (\textit{fail-fast}).
    \item \textbf{Valores por defecto seguros:} Variables no críticas como SMTP se manejan como \texttt{Optional} o \texttt{str | None}.
\end{enumerate}

\subsection{Propiedades de \texttt{model\_config}}
El diccionario de configuración dicta el comportamiento de carga para Pydantic:
\begin{itemize}
    \item \texttt{env\_file=".env"}: Indica que las variables pueden leerse de un archivo local.
    \item \texttt{env\_file\_encoding="utf-8"}: Soporta caracteres especiales en credenciales.
    \item \texttt{case\_sensitive=True}: Diferencia entre MAYÚSCULAS y minúsculas para los nombres de las variables.
    \item \texttt{extra="ignore"}: Ignora variables presentes en el archivo \texttt{.env} que no estén definidas explícitamente en la clase.
\end{itemize}

\subsection{Código Fuente de \texttt{config.py}}

\begin{lstlisting}[language=Python, caption=app/core/config.py]
"""
Proyecto: Plataforma de Asistencia Biometrica Escolar
Modulo: Infraestructura Base
Archivo: config.py
Descripcion: Centraliza y valida variables de entorno mediante pydantic.
Autor: Marco Antonio Pena Vargas
Fecha de Creacion: 6 Agosto 2026
"""

from pydantic_settings import BaseSettings, SettingsConfigDict

class Settings(BaseSettings):
    """
    Configuracion global de la aplicacion validada por Pydantic.
    """
    # Configuracion general
    Nombre_Proyecto: str = "Plataforma Escolar con validacion Biometrica"
    Version_Proyecto: str = "1.0.0"
    API_Version: str = "/api/v1"
    ENV: str = "development" 
    DEBUG: bool = False

    # Configuracion de base de datos
    DB_URL: str
    SUPABASE_URL: str
    SUPABASE_ANON_KEY: str
    SUPABASE_SERVICE_ROLE_KEY: str | None = None

    # Autenticacion JWT
    JWT_SECRET: str
    JWT_ALGORITHM: str = "HS256"
    ACCESS_TOKEN_EXPIRE_MINUTES: int = 30

    # Archivos
    UPLOAD_DIRECTORY: str = "uploads"
    CLOUDINARY_URL: str | None = None

    # SMTP Correo
    SMTP_SERVER: str | None = None
    SMTP_PORT: int = 587
    SMTP_USERNAME: str | None = None
    SMTP_PASSWORD: str | None = None
    SMTP_USE_TLS: bool = True

    model_config = SettingsConfigDict(
        env_file=".env",
        env_file_encoding="utf-8",
        case_sensitive=True,
        extra="ignore",
    )

# Instancia unica de configuracion global declarada FUERA de la clase
settings = Settings()
\end{lstlisting}

% --- 5. VARIABLES DE ENTORNO ---
\section{Variables de Entorno y Seguridad}

\subsection{Diferencia Crítica de Permisos en Supabase}
El sistema interactúa con PostgreSQL utilizando dos niveles de credenciales que \textbf{no deben confundirse}:
\begin{itemize}
    \item \textbf{\texttt{SUPABASE\_ANON\_KEY}:} Es la clave pública de la API, restringida por las políticas de seguridad (RLS).
    \item \textbf{\texttt{SUPABASE\_SERVICE\_ROLE\_KEY}:} Es una llave maestra (\textit{admin key}) que ignora las reglas RLS. \textbf{Exclusiva para el Backend.} Jamás debe exponerse al frontend.
\end{itemize}

\subsection{Archivo \texttt{.env} de Desarrollo}
Ubicado en \texttt{backend/.env} (Ignorado en \texttt{.gitignore}). Define los valores reales:

\begin{lstlisting}[language=bash, caption=backend/.env]
Nombre_Proyecto=Plataforma Escolar con validacion Biometrica
Version_Proyecto=1.0.0
API_Version=/api/v1
ENV=development
DEBUG=true

DB_URL=TU_DATABASE_URL
SUPABASE_URL=TU_SUPABASE_URL
SUPABASE_ANON_KEY=TU_SUPABASE_ANON_KEY
SUPABASE_SERVICE_ROLE_KEY=TU_SERVICE_ROLE_KEY

JWT_SECRET=TU_SECRETO_JWT
JWT_ALGORITHM=HS256
ACCESS_TOKEN_EXPIRE_MINUTES=30

UPLOAD_DIRECTORY=uploads
CLOUDINARY_URL=

SMTP_SERVER=
SMTP_PORT=587
SMTP_USER=
SMTP_PASSWORD=
SMTP_USE_TLS=true
\end{lstlisting}

% --- 6. TROUBLESHOOTING ---
\section{Registro de Errores y Soluciones (Troubleshooting)}

Durante la configuración inicial del módulo \texttt{config.py} y la carga de variables mediante \texttt{pydantic-settings}, se documentaron y resolvieron las siguientes incidencias para asegurar la estabilidad del entorno de desarrollo:

\subsection{Incidencia 1: Error de Ámbito en la Instanciación de la Clase}
\begin{itemize}
    \item \textbf{Problema:} \texttt{NameError: name 'Settings' is not defined} al ejecutar las pruebas en \texttt{test\_config.py}.
    \item \textbf{Causa:} Se definió la instancia global \texttt{settings = Settings()} dentro del bloque de la propia clase \texttt{Settings}. En Python, una clase no está completamente definida hasta que se ejecuta todo su bloque; por lo tanto, no se puede referenciar el nombre de la clase dentro de su propio cuerpo.
    \item \textbf{Solución aplicada:} Se trasladó la creación de la instancia \texttt{settings = Settings()} \textbf{fuera de la definición de la clase}, dejándola al final del archivo como una variable global en el módulo \texttt{config.py}.
    \item \textbf{Impacto:} La clase se construye exitosamente y la configuración puede ser importada por otras dependencias del sistema.
    \item \textbf{Archivos modificados:} \texttt{app/core/config.py}.
\end{itemize}

\subsection{Incidencia 2: Error de Validación Pydantic por Nomenclatura}
\begin{itemize}
    \item \textbf{Problema:} \texttt{ValidationError: 1 validation error for Settings. DB\_URL Field required [type=missing]}.
    \item \textbf{Causa:} El archivo \texttt{.env} contenía la variable declarada con un error tipográfico (\texttt{DB\_UR}), mientras que el modelo \texttt{Settings} en el código requería estrictamente \texttt{DB\_URL}. Al no existir una coincidencia exacta de nombres, Pydantic bloqueó de inmediato el inicio del servidor, aplicando correctamente el principio \textit{fail-fast}.
    \item \textbf{Solución aplicada:} Se renombró la variable en el archivo \texttt{.env} de \texttt{DB\_UR} a \texttt{DB\_URL}. Adicionalmente, se verificó que el resto de las variables (como \texttt{JWT\_EXPIRE\_MINUTES} a \texttt{ACCESS\_TOKEN\_EXPIRE\_MINUTES}) coincidieran perfectamente entre el archivo físico y la clase de validación.
    \item \textbf{Impacto:} Eliminación del error de validación, logrando la lectura íntegra y segura de todas las credenciales de base de datos.
    \item \textbf{Archivos modificados:} \texttt{.env}.
\end{itemize}

\end{document}